\section{Preparation-independent critical-layer geometry}
\label{sec:unprepared-outer-skeleton}

The finite-section noncanonicity result raises a sharper question: which
outer features are nevertheless fixed by the physical equation?  The
following result answers this at the critical-layer level and shows precisely
why it does not select a complete history.  The splitting is for the leading
frozen-resource fast-voltage matrix obtained by omitting the
\(O(\delta^2)\) current and delayed variational terms; the critical curve is
not a finite-\(\delta\) slow solution.

Put
\begin{equation}
 \mathcal A_N=D(P_N-\Id),\qquad
 \mathcal N_N(x)=c_N\circ x^{\circ2}
              +\frac\beta3x^{\circ3}.                 \label{eq:outer-fast-nonlinearity}
\end{equation}

\begin{proposition}[Uniform critical curve and leading frozen-voltage splitting]
\label[proposition]{prop:unprepared-outer-skeleton}
Assume the common Markov, curvature, and operator bounds in
\cref{eq:shared-markov,eq:shared-curvature,eq:shared-layer-bounds}.
There are \(r_0,C,C_P>0\), independent of \(N\) and of the delay layers,
with the following properties.

For \(|r|\le r_0\), there is a solution
\(z_{0,N}(r)\in E_N\) of
\begin{equation}
 A_Nz_{0,N}(r)
 =P_{\perp,N}\mathcal N_N
       \bigl(r\mathbf1+z_{0,N}(r)\bigr),
 \qquad \|z_{0,N}(r)\|_N\le Cr^2,                  \label{eq:outer-critical-transverse}
\end{equation}
and it is the unique solution in the displayed \(Cr^2\)-ball.
Define
\begin{align}
 x_{0,N}(r)&=r\mathbf1+z_{0,N}(r),\notag\\
 V_{0,N}(r)&=\mathbf1+x_{0,N}(r),\notag\\
 w_{0,N}(r)&=\frac23-\pi_N^T\mathcal N_N(x_{0,N}(r)).
                                                               \label{eq:outer-critical-curve}
\end{align}
Then the constant history with value \(V_{0,N}(r)\), at frozen resource
\(w_{0,N}(r)\), is an exact zero of the unprepared voltage field for every
positive \(\delta\), every admitted redistribution, and either recovery law.  In
particular, this critical curve uses no cutoff or completion.  Uniformly for
\(0\le j\le3\),
\begin{align}
 z_{0,N}(r)
 &=g_Nr^2+g_{3,N}r^3+R^z_N(r),
 &\|\partial_r^jR^z_N(r)\|_N&\le C|r|^{4-j},\notag\\
 w_{0,N}(r)
 &=\frac23-\alpha r^2-\kappa_Nr^3+R^w_N(r),
 &|\partial_r^jR^w_N(r)|&\le C|r|^{4-j},             \label{eq:outer-critical-expansion}
\end{align}
where
\begin{equation}
 g_{3,N}=2A_N^{-1}P_{\perp,N}(c_N\circ g_N).         \label{eq:outer-g3}
\end{equation}
The two curve components are \(C^\infty\) in \(r\).  More precisely, for
every fixed finite \(s\),
\begin{equation}
 \max_{0\le j\le s}\sup_{|r|\le r_0}
 \left\{\|\partial_r^jz_{0,N}(r)\|_N
             +|\partial_r^jw_{0,N}(r)|\right\}
 \le C_s,                                           \label{eq:outer-critical-high-regularity}
\end{equation}
where \(C_s\) is independent of \(N\).

The leading frozen-resource fast-voltage matrix along this curve is
\begin{equation}
 \mathcal J^{\rm fv}_{N}(r)=\mathcal A_N-\diag\!\left(
       2c_N\circ x_{0,N}(r)+\beta x_{0,N}(r)^{\circ2}
                                  \right).            \label{eq:outer-layer-jacobian}
\end{equation}
It has a real simple eigenvalue \(\lambda_{c,N}(r)\), continued from the
collective zero eigenvalue at the fold, and corresponding right eigenvector
and left eigenfunctional \(e_{c,N}(r)\), \(\ell_{c,N}(r)\), normalized by
\begin{equation}
 \ell_{c,N}(r)\mathbf1=1,
 \qquad \ell_{c,N}(r)e_{c,N}(r)=1.                  \label{eq:outer-eigen-normalization}
\end{equation}
Writing \(O_{\rm unif}\) for a remainder whose bound is independent of
\(N\), they obey, in the vector and dual norms,
\begin{align}
 e_{c,N}(r)&=\mathbf1+2rg_N+O_{\rm unif}(r^2),\notag\\
 \ell_{c,N}(r)&=\pi_N^T
   +2r\,\pi_N^T\diag(c_N)A_N^{-1}P_{\perp,N}+O_{\rm unif}(r^2),\notag\\
 \lambda_{c,N}(r)&=-2\alpha r-3\kappa_Nr^2+O_{\rm unif}(r^3),
                                                               \label{eq:outer-eigen-expansion}
\end{align}
where all remainder constants are independent of \(N\).  If
\(P^s_N(r)=\Id-e_{c,N}(r)\ell_{c,N}(r)\), then
\begin{equation}
 \|e^{\mathcal J^{\rm fv}_{N}(r)t}P^s_N(r)\|
 \le C_Pe^{-D\gamma t/2},\qquad t\ge0.              \label{eq:outer-stable-layer-semigroup}
\end{equation}
Consequently the leading frozen-resource voltage layer is attracting for
\(0<r\le r_0\), whereas for \(-r_0\le r<0\) it has exactly one unstable
direction.  This statement excludes both the resource zero direction and all
positive-\(\delta\) current and delayed variational terms.

There is also an exact sign identity.  With
\begin{equation}
 \Theta_N(r)=\ell_{c,N}(r)V_{0,N}'(r),               \label{eq:outer-theta}
\end{equation}
one has \(\Theta_N(r)=1+O_{\rm unif}(r^2)\); after decreasing \(r_0\),
\begin{equation}
 \frac12\le\Theta_N(r)\le\frac32,
 \qquad
 \lambda_{c,N}(r)=\frac{w_{0,N}'(r)}{\Theta_N(r)}.  \label{eq:outer-eigen-identity}
\end{equation}
\end{proposition}

\begin{corollary}[Truncated action for the reduced critical-layer equation]
\label[corollary]{cor:outer-reduced-actions}
Retain \cref{prop:unprepared-outer-skeleton}.
For the shared-resource and sensed-recovery laws, respectively, set
\begin{align}
 \sigma_N^{\rm sh}(r)&=r,\notag\\
 \sigma_N^{\rm rec}(r)&=\pi_N^T\!\left[
   \varpi_N\circ x_{0,N}(r)
   +d_N^{\rm rec}\circ x_{0,N}(r)^{\circ2}\right]. \label{eq:outer-recovery-source}
\end{align}
For the second formula, assume additionally \cref{eq:sensing-bounds}.  Then
\begin{equation}
 \sigma_N^{\rm rec}(r)=r+\chi_N^{\rm rec}r^2+O_{\rm unif}(r^3).
                                                               \label{eq:outer-recovery-expansion}
\end{equation}
Fix \(p>4\) and \(\nu_*>0\), restrict to \(|\nu|\le\nu_*\), and let
\begin{equation}
 S_\delta=\sqrt{2p\log(1/\delta)},\qquad
 \rho_\delta=\frac{\delta S_\delta}{2\alpha},       \label{eq:outer-matching-radius}
\end{equation}
and use slow time \(T=\delta^2t\).  There are fixed
\(0<r_{\rm out}<r_0\) and \(\delta_0>0\), depending only on
\(p,\nu_*\) and the common bounds, such that, for either
\(\diamond\in\{\mathrm{sh},\mathrm{rec}\}\) and every
\(0<\delta<\delta_0\),
\begin{equation}
 q^\diamond_{N,\delta,\nu}(r)
 :=\frac{\sigma_N^\diamond(r)-\delta^2\nu}{w_{0,N}'(r)}
 =-\frac1{2\alpha}
   +O_{\rm unif}\!\left(|r|+\frac{\delta^2}{|r|}\right)       \label{eq:outer-reduced-speed}
\end{equation}
on \(\rho_\delta\le|r|\le r_{\rm out}\), and
\begin{equation}
 -\frac{3}{2\alpha}\le q^\diamond_{N,\delta,\nu}(r)
 \le-\frac1{6\alpha}.                               \label{eq:outer-speed-sign}
\end{equation}
The two truncated reduced actions
\begin{align}
 \mathcal A^a_{N,\diamond}
 &=\int_{\rho_\delta}^{r_{\rm out}}
      \frac{\lambda_{c,N}(r)}
           {q^\diamond_{N,\delta,\nu}(r)}\,\dd r,\notag\\
 \mathcal A^r_{N,\diamond}
 &=-\int_{-r_{\rm out}}^{-\rho_\delta}
      \frac{\lambda_{c,N}(r)}
           {q^\diamond_{N,\delta,\nu}(r)}\,\dd r       \label{eq:outer-normal-actions}
\end{align}
are positive for \(0<\delta<\delta_0\).  These are truncated actions of the
formal reduced critical-curve equation, not actions of finite-\(\delta\)
invariant RFDE histories.  If \(\rho_\delta\le r_{\rm out}/2\), then, uniformly in
\(N\), \(\nu\), and the two recovery laws,
\begin{equation}
 \min\{\mathcal A^a_{N,\diamond},
        \mathcal A^r_{N,\diamond}\}
 \ge \frac{\alpha^2r_{\rm out}^2}{6}=:\mathcal A_*>0,
                                                               \label{eq:outer-action-lower-bound}
\end{equation}
and
\begin{equation}
 \mathcal A^{a/r}_{N,\diamond}
 =2\alpha^2(r_{\rm out}^2-\rho_\delta^2)
   +O_{\rm unif}(r_{\rm out}^3+\delta^2r_{\rm out}). \label{eq:outer-action-expansion}
\end{equation}
\end{corollary}

The preceding matrix splitting does have a positive-\(\delta\),
voltage-history-space consequence, provided the resource coordinate is frozen.
This is the first ambient RFDE estimate needed by an outer-history
construction.  It is convenient to use fold time \(s=\delta t\), for which
the history length is fixed rather than \(O(\delta^{-1})\).

For
\(\eta=(E_0,\ldots,E_m)\in\mathfrak T_N^{\rm red}\), write
\(L_{k,N}(\eta)=B_{k,N}+E_k\), and put
\begin{equation}
 \mathscr X_N=C([-\theta_m,0];\R^N),\qquad
 \mathfrak a_N(r)=
 \left\{\frac1{|\lambda_{c,N}(r)|}
             +\frac2{D\gamma}\right\}^{-1}.         \label{eq:frozen-history-a}
\end{equation}
Fix \(0<\vartheta<1/4\).  For \(0<\delta<e^{-1}\) and \(r\ne0\), define
the smooth capped rate and its associated fading-history norm by
\begin{align}
 \kappa_{N,\delta}(r)
 &=\vartheta\left\{
       \frac{\delta}{\mathfrak a_N(r)}
       +\frac{\theta_m}{\log(1/\delta)}\right\}^{-1},
                                                        \label{eq:frozen-history-rate}\\
 \|\phi\|_{N,r,\delta}^{\sharp}
 &=\sup_{-\theta_m\le\theta\le0}
       e^{2\kappa_{N,\delta}(r)\theta}\|\phi(\theta)\|_N.
                                                        \label{eq:frozen-history-norm}
\end{align}
The harmless case \(\theta_m=0\) is interpreted by deleting the second
term in braces and reducing the history equation to an ODE.

\begin{proposition}[Frozen-resource voltage-history splitting]
\label[proposition]{prop:frozen-voltage-history-splitting}
Retain \cref{prop:unprepared-outer-skeleton,cor:outer-reduced-actions} and
the structural-ball layers \cref{eq:structural-ball-layers}.  After
decreasing \(r_{\rm out},\delta_0,\eta_0>0\), independently of \(N\), the
following holds whenever
\[
 0<\delta<\delta_0,\qquad
 \rho_\delta\le |r|\le r_{\rm out},\qquad
 \|\eta\|_{\mathfrak T_N}<\eta_0.
\]
On \(\mathscr X_N\), consider the frozen-resource voltage RFDE
\begin{equation}
 u'(s)=\delta^{-1}\mathcal J_N^{\rm fv}(r)u(s)
 +\delta K\sum_{k=0}^mL_{k,N}(\eta)
       \{u(s)-u(s-\theta_k)\}.                      \label{eq:frozen-voltage-rfde}
\end{equation}
Let \(\mathcal T_{N,\delta,r,\eta}(s)\) denote its solution semigroup.

If \(r>0\), the whole history space is exponentially stable and
\begin{equation}
 \|\mathcal T_{N,\delta,r,\eta}(s)\phi\|_{N,r,\delta}^{\sharp}
 \le C e^{-\kappa_{N,\delta}(r)s}
       \|\phi\|_{N,r,\delta}^{\sharp},qquad s\ge0.  \label{eq:frozen-attracting-history-estimate}
\end{equation}
If \(r<0\), there is an invariant splitting
\begin{equation}
 \mathscr X_N=\mathscr W^s_{N,\delta,r,\eta}
                 \oplus\mathscr W^u_{N,\delta,r,\eta},
 \qquad \operatorname{codim}\mathscr W^s=1,
 \qquad \dim\mathscr W^u=1,                         \label{eq:frozen-repelling-history-splitting}
\end{equation}
whose projections have norm at most \(C\) in
\cref{eq:frozen-history-norm}.  On the stable space,
\begin{equation}
 \|\mathcal T(s)\phi^s\|_{N,r,\delta}^{\sharp}
 \le C e^{-\kappa_{N,\delta}(r)s}
       \|\phi^s\|_{N,r,\delta}^{\sharp},\qquad s\ge0.          \label{eq:frozen-repelling-stable-estimate}
\end{equation}
The unstable space is spanned by
\begin{equation}
 h^u_{N,\delta,r,\eta}(\theta)
 =e^{\lambda^u_{N,\delta,r,\eta}\theta}
       e^u_{N,\delta,r,\eta},qquad -\theta_m\le\theta\le0, \label{eq:frozen-unstable-history}
\end{equation}
where \(\ell_{c,N}(r)e^u_{N,\delta,r,\eta}=1\) and
\begin{equation}
 \left|\lambda^u_{N,\delta,r,\eta}
       -\delta^{-1}\lambda_{c,N}(r)\right|\le C\delta,
 \qquad
 \|e^u_{N,\delta,r,\eta}-e_{c,N}(r)\|_N\le C\delta^2,       \label{eq:frozen-unstable-root-estimate}
\end{equation}
and the inverse restricted to this line satisfies
\begin{equation}
 \|\mathcal T(-s)|_{\mathscr W^u}\|
 \le C e^{-\lambda^u_{N,\delta,r,\eta}s},\qquad s\ge0.     \label{eq:frozen-unstable-inverse-estimate}
\end{equation}
Here \(\mathcal T(-s)|_{\mathscr W^u}\) is only the inverse on the
one-dimensional invariant line; no backward RFDE semiflow is asserted.

The bundles, projections, unstable root, and normalized unstable history are
twice Fr\'{e}chet differentiable in \(\eta\), in fixed graph charts, with
common operator-norm bounds through order two.  The proof is controlled by
the dimension-free quantity
\begin{equation}
 \epsilon_{N,\delta}(r)
 :=C\frac{\delta^2
       \{1+e^{2\kappa_{N,\delta}(r)\theta_m}\}}
          {\mathfrak a_N(r)}
 \le C\frac{\delta^{1-2\vartheta}}{S_\delta}=o(1).  \label{eq:frozen-history-smallness}
\end{equation}
In particular, the attracting equation has no characteristic root with
real part greater than \(-\kappa_{N,\delta}(r)\); on the repelling branch the
only such root is the real simple root in
\cref{eq:frozen-unstable-root-estimate}.
\end{proposition}

The fading factor \(2\kappa\) in the phase norm and the smaller decay rate
\(\kappa\) in the conclusion are intentional.  The margin controls the
translation of prescribed old history before one full delay has elapsed;
using the same exponent in both places would leave no perturbative gap.
For clarity, this is a fading-norm estimate: its comparison with the ordinary
supremum norm is
\begin{equation}
 \|\phi\|^\sharp_{N,r,\delta}\le\|\phi\|_\infty
 \le e^{2\kappa_{N,\delta}(r)\theta_m}
       \|\phi\|^\sharp_{N,r,\delta}
 \le\delta^{-2\vartheta}\|\phi\|^\sharp_{N,r,\delta}.         \label{eq:frozen-history-norm-equivalence}
\end{equation}
The cap is also necessary at the level of this model class: long-delay weak
modes need not retain a \(\delta\)-independent fast-time decay rate.

\begin{remark}[Why the resource direction must be phase-quotiented]
\label[remark]{rem:frozen-resource-phase-quotient}
\Cref{prop:frozen-voltage-history-splitting} cannot be upgraded by simply
adjoining the frozen resource equation and retaining the phrase ``one
unstable direction.''  In the admissible homogeneous, delay-free collective
subclass \(c_N=\alpha\mathbf1\), the collective voltage--resource block in
the original time is
\[
 \dot u=\lambda_{c,N}(r)u-\omega,
 \qquad \dot\omega=\delta^2u,
\]
and hence has characteristic polynomial
\begin{equation}
 \lambda^2-\lambda_{c,N}(r)\lambda+\delta^2.         \label{eq:full-frozen-resource-polynomial}
\end{equation}
For \(r<0\) with \(|r|\ge\rho_\delta\), both roots are real and positive
for all sufficiently small \(\delta\).  Thus the unquotiented full frozen
system has two unstable collective--resource roots on the repelling branch.
The additional slow root signals a degree of freedom that can be related to
the longitudinal bundle only after a genuine slow history has been
constructed.  A physical outer theorem must then prove either a
stable--longitudinal--strong-unstable trichotomy or, equivalently, a
stable/one-strong-unstable dichotomy for the correctly gauged quotient
process.  The raw full frozen system has neither of those already-identified
structures.
\end{remark}

The true outer-history problem can be written without a completion or
cutoff.  The next proposition records the exact equation, including the
state-dependent backtrack generated by the physical slow speed.
On the fixed fold-time phase space put
\begin{equation}
 \mathscr Y_N=C([-\theta_m,0];\R^N)\times\R,
 \qquad
 \varrho_N(\phi,\omega)=\pi_N^T\phi(0).             \label{eq:outer-stationary-gauge-row}
\end{equation}

\begin{proposition}[Exact uncut outer-history equation]
\label[proposition]{prop:exact-outer-history-equation}
Fix \(N,\delta,\nu\), restrict to the shared-resource law, take
\(\eta\in\mathfrak T_N^{\rm red}\), use
\(L_{k,N}(\eta)=B_{k,N}+E_k\), and let \(J\subset\R\) be an interval.
Suppose \(q\in C^1(J)\) is nonzero and its scalar flow
\(\Phi_q^T\) is defined for the required backward times.  Put
\begin{equation}
 r_k(r)=\Phi_q^{-\delta\theta_k}(r),                 \label{eq:exact-outer-backtrack}
\end{equation}
and assume \(r_k(r)\in J\) for every retained \(r\) and \(k\).  Let
\[
 V(r)=\mathbf1+r\mathbf1+z(r),\qquad
 z(r)\in E_N,\qquad w(r)\in\R,
\]
with \(z,w\in C^1(J)\).  Then the following are equivalent:

\begin{enumerate}[label=\textup{(\roman*)},leftmargin=2em]
\item for every retained \(r_0\), these functions generate a physical
solution of the unmodified structural-ball RFDE by
\begin{equation}
 r(T)=\Phi_q^T(r_0),\qquad
 \bar v(t)=V(r(\delta^2t)),\qquad
 \bar w(t)=w(r(\delta^2t))                          \label{eq:exact-outer-generated-solution}
\end{equation}
on every orbit segment for which the orbit and all delayed backtracks remain
in \(J\);
\item the identities
\begin{align}
 \delta^2q(r)\{\mathbf1+z'(r)\}
 ={}&\{2/3-w(r)\}\mathbf1
 -\mathcal N_N(r\mathbf1+z(r))+A_Nz(r)\notag\\
 &+\delta^2K\sum_{k=0}^mL_{k,N}(\eta)
 \bigl[\{r-r_k(r)\}\mathbf1+z(r)-z(r_k(r))\bigr],
                                                        \label{eq:exact-outer-voltage-invariance}\\
 q(r)w'(r)={}&r-\delta^2\nu.                         \label{eq:exact-outer-resource-invariance}
\end{align}
hold at every retained \(r\).
\end{enumerate}
Every \(C^1\) solution of these identities is automatically \(C^2\) on the
interior of the retained set.

Equivalently, the full fold-time history over the curve is
\begin{equation}
 \mathcal H_{V,w,q}(r)
 =\left(
   \theta\longmapsto V(\Phi_q^{\delta\theta}(r)),
   \;w(r)\right)\in\mathscr Y_N.                    \label{eq:exact-outer-history-embedding}
\end{equation}
It satisfies
\begin{equation}
 \varrho_N\partial_r\mathcal H_{V,w,q}(r)=1.        \label{eq:exact-outer-history-phase}
\end{equation}
Along a generated orbit define \(r_*(s)=r(\delta s)\).  Then
\(r_*'(s)=\delta q(r_*(s))\), and the normalized
time-tangent history in
\cref{eq:physical-normalized-tangent} is exactly
\begin{equation}
 \tau_*(s)=
 \partial_r\mathcal H_{V,w,q}(r_*(s)).              \label{eq:exact-outer-history-tangent}
\end{equation}
\end{proposition}

The state-dependent backtrack in
\cref{eq:exact-outer-backtrack} has to be differentiated in the tracker
problem.  The next estimate makes explicit the only relative loss: the delay
shift divided by the distance from the fold.

For an interval \(J\) contained in one component of \(\R\setminus\{0\}\),
put
\begin{equation}
 \|f\|_{k,\rho;J}
 =\max_{0\le j\le k}\sup_{r\in J}
       |r|^j\|\partial_r^jf(r)\|.                  \label{eq:backtrack-weighted-norm}
\end{equation}
The definition applies to scalar-, vector-, and operator-valued functions.

\begin{proposition}[Weighted calculus for the physical backtrack]
\label[proposition]{prop:physical-backtrack-calculus}
Fix an integer \(k\ge2\), \(0<q_-<q_+\), and \(M>0\).  Let
\(I\Subset J\) be compact intervals in one component of
\(\R\setminus\{0\}\) such that
\begin{equation}
 |r|\ge\rho\quad(r\in I),
 \qquad |r|\ge\rho/2\quad(r\in J).                \label{eq:backtrack-collar-geometry}
\end{equation}
Suppose \(q\in C^{k+2}(J)\) satisfies
\begin{equation}
 q_-\le-q(r)\le q_+,
 \qquad \|q\|_{k+2,\rho;J}\le M,                 \label{eq:backtrack-q-class}
\end{equation}
Suppose moreover that one fixed interval \(J_0\Subset J\) contains the whole
backward flow tube
\begin{equation}
 \left\{\Phi_q^{-\tau}(r):
       r\in\overline I,\ 0\le\tau\le\delta\theta_m\right\}
 \subset J_0.                                      \label{eq:backtrack-flow-tube}
\end{equation}
Put \(r_j(r)=\Phi_q^{-\delta\theta_j}(r)\).  The collars and \(\rho\) may
vary with \(\delta\).  There are \(\epsilon_0,C>0\), depending only on
\(k,\theta_m,q_-,q_+\), and \(M\), but not on \(N,\delta,\rho\), or the
locations and lengths of the displayed intervals, such that, whenever
\begin{equation}
 \epsilon_\rho:=\frac{\delta\theta_m}{\rho}
 \le\epsilon_0,                                    \label{eq:backtrack-relative-shift}
\end{equation}
the following estimates hold for every \(j\):
\begin{align}
 |r_j(r)-r|&\le q_+\delta\theta_j,
 &r_j'(r)&=\frac{q(r_j(r))}{q(r)},                  \label{eq:backtrack-basic-identities}\\
 \|r_j-\Id\|_{k,\rho;I}&\le C\delta\theta_j.    \label{eq:backtrack-map-bound}
\end{align}
On the open subset of \(C^{k+2}(J)\), endowed with
\cref{eq:backtrack-weighted-norm}, whose swept flow tube remains in
\(\operatorname{int}J\), the map
\[
 q\longmapsto r_j-\Id:
 C^{k+2}(J)\longrightarrow C^k(I)
\]
is three times Fr\'{e}chet differentiable.  For
\(h,h_1,h_2,h_3\in C^{k+2}(J)\),
\begin{align}
 D_qr_j[h]
 &=q(r_j)\int_r^{r_j}\frac{h(u)}{q(u)^2}\,\dd u,  \label{eq:backtrack-first-variation}\\
 \|D_qr_j[h]\|_{k,\rho;I}
 &\le C\delta\theta_j\|h\|_{k+2,\rho;J},\notag\\
 \|D_q^2r_j[h_1,h_2]\|_{k,\rho;I}
 &\le C\delta\theta_j
       \|h_1\|_{k+2,\rho;J}\|h_2\|_{k+2,\rho;J},\notag\\
 \|D_q^3r_j[h_1,h_2,h_3]\|_{k,\rho;I}
 &\le C\delta\theta_j
       \prod_{a=1}^3\|h_a\|_{k+2,\rho;J}.
                                                               \label{eq:backtrack-frechet-bounds}
\end{align}
If \(f\in C^{k+3}(J;X)\), where \(X\) is any Banach space, then
\begin{align}
 \|f\circ r_j\|_{k,\rho;I}
 &\le C\|f\|_{k,\rho;J},                           \label{eq:backtrack-composition-bound}\\
 \|f\circ r_j-f\|_{k,\rho;I}
 &\le C\epsilon_\rho\|f\|_{k+1,\rho;J},          \label{eq:backtrack-composition-difference}\\
 \|D_q(f\circ r_j)[h]\|_{k,\rho;I}
 &\le C\epsilon_\rho\|f\|_{k+1,\rho;J}
                    \|h\|_{k+2,\rho;J},\notag\\
 \|D_q^2(f\circ r_j)[h_1,h_2]\|_{k,\rho;I}
 &\le C\epsilon_\rho\|f\|_{k+2,\rho;J}
       \|h_1\|_{k+2,\rho;J}\|h_2\|_{k+2,\rho;J},\notag\\
 \|D_q^3(f\circ r_j)[h_1,h_2,h_3]\|_{k,\rho;I}
 &\le C\epsilon_\rho\|f\|_{k+3,\rho;J}
       \prod_{a=1}^3\|h_a\|_{k+2,\rho;J}.
                                                               \label{eq:backtrack-composition-frechet}
\end{align}
More precisely, on any product neighborhood with one common flow tube, put
\[
 \mathcal C_j(f,q)=f\circ r_j(q),\qquad
 \Delta_j(f,q)=f\circ r_j(q)-f|_I.
\]
Both maps from
\(C^{k+3}(J;X)\times C^{k+2}(J)\) to \(C^k(I;X)\) are three times
Fr\'{e}chet differentiable.  If \(1\le b\le3\) and
\((g_a,h_a)\) are product-space directions, then
\begin{align}
 &\|D^b\Delta_j(f,q)
       [(g_1,h_1),\ldots,(g_b,h_b)]\|_{k,\rho;I}\notag\\
 &\quad\le C\epsilon_\rho\left\{
      \|f\|_{k+3,\rho;J}\prod_{a=1}^b\|h_a\|_{k+2,\rho;J}
      +\sum_{\ell=1}^b\|g_\ell\|_{k+3,\rho;J}
          \prod_{a\ne\ell}\|h_a\|_{k+2,\rho;J}\right\}.
                                                               \label{eq:backtrack-joint-difference}
\end{align}
The same estimate without the factor \(\epsilon_\rho\) holds for
\(D^b\mathcal C_j\).  Thus a parameter-dependent pair \((f,q)\), with a
common flow tube and common \(C^1_\nu C^2_\eta\) bounds in these two Banach
spaces, produces delayed compositions and delayed differences with the same
rectangular parameter regularity, including the mixed
\(\partial_\nu\partial_{\eta\eta}\) derivative.  The same conclusion holds
for every rectangular jet of total order at most three, in particular also
for \(C^2_\nu C^1_\eta\) and the mixed
\(\partial_{\nu\nu}\partial_\eta\) derivative.  A raw backtrack variation has
the absolute factor \(C\delta\theta_j\) in
\cref{eq:backtrack-frechet-bounds}; once it enters a weighted delayed
composition, its relative factor is \(C\epsilon_\rho\).  At
\(\rho=\rho_\delta\), this factor is
\begin{equation}
 \epsilon_{\rho_\delta}\le\frac{C}{S_\delta}.       \label{eq:backtrack-logarithmic-loss}
\end{equation}
All constants in these conclusions are uniform in \(N,\delta\), and the
admissible collar family just specified.
\end{proposition}

The proposition is an exact reduction, not an existence assertion.
In particular, imposing boundary data on
\cref{eq:exact-outer-voltage-invariance,eq:exact-outer-resource-invariance}
can still select different representatives.  The missing theorem must
construct a nonempty physical representative class and control that
dependence in full-history parameter norms.

The remaining principal tracker linearization has an apparent derivative
coupling from the transverse curve into the slow speed.  The next result
removes that singularity exactly and proves the endpoint Green splitting
that is available before compatible complete collars have been constructed.
It is restricted to the shared-resource law.  Put
\begin{align}
 B_N(r)
 &=A_N-P_{\perp,N}D\mathcal N_N(x_{0,N}(r))|_{E_N},\notag\\
 \ell_N(r)
 &=\pi_N^TD\mathcal N_N(x_{0,N}(r))|_{E_N}.        \label{eq:principal-tracker-coefficients}
\end{align}

\begin{proposition}[Endpoint Green splitting for the principal tracker]
\label[proposition]{prop:principal-tracker-endpoint-green}
Retain \cref{prop:unprepared-outer-skeleton,cor:outer-reduced-actions}, use
the shared-resource speed \(q_0=q^{\rm sh}_{N,\delta,\nu}\), fix an integer
regularity order \(s\ge1\), and put
\begin{equation}
 \varepsilon=\delta^2,\qquad
 I_\delta^a=[\rho_\delta,r_{\rm out}],\qquad
 I_\delta^r=[-r_{\rm out},-\rho_\delta].           \label{eq:principal-green-intervals}
\end{equation}
After decreasing \(r_{\rm out}\) and then \(\delta_0\), independently of
\(N\), the following holds for \(|\nu|\le\nu_*\) and
\(0<\delta<\delta_0\).  On either interval,
\begin{align}
 -q_+\le q_0(r)\le-q_-<0,\qquad
 &c_w|r|\le|w_{0,N}'(r)|\le C_w|r|,\qquad
 rw_{0,N}'(r)<0,                                   \label{eq:principal-green-signs}\\
 \max_{0\le j\le s+1}|r|^j
 \bigl\{\|\partial_r^j(B_N-A_N)\|
 &+\|\partial_r^jz_{0,N}'\|_N
  +\|\partial_r^j\ell_N\|
  +|\partial_r^jw_{0,N}'|\bigr\}
 \le C_s|r|,                                      \label{eq:principal-green-coefficients}\\
 \max_{0\le j\le s+1}|r|^j
 |\partial_r^jq_0(r)|&\le C_s.                    \label{eq:principal-green-q-bounds}
\end{align}
All displayed constants are dimension uniform.

For
\begin{equation}
 \mathcal L_{N,\delta}\binom ZQ
 =\binom{
   \varepsilon q_0Z'-B_NZ+\varepsilon z_{0,N}'Q}
  {-\varepsilon q_0Q'+w_{0,N}'Q-q_0(\ell_NZ)'}
 =\binom fg,                                       \label{eq:principal-tracker-operator}
\end{equation}
set
\begin{equation}
 \mathfrak p=\varepsilon Q+\ell_NZ.                \label{eq:principal-green-p-coordinate}
\end{equation}
For \(I\in\{I_\delta^a,I_\delta^r\}\), define
\begin{align}
 \|Z\|_{\mathsf Z_\varepsilon^s(I)}
 &=\varepsilon^{-1}\max_{0\le j\le s}\sup_{r\in I}
       |r|^j\|\partial_r^jZ(r)\|_N,\notag\\
 \|\mathfrak p\|_{\mathsf P_\varepsilon^s(I)}
 &=\varepsilon^{-1}\max_{0\le j\le s}\sup_{r\in I}
       |r|^{j-1}|\partial_r^j\mathfrak p(r)|,\notag\\
 \|Q\|_{\mathsf Q_\varepsilon^j(I)}
 &=\varepsilon^{-1}\max_{0\le i\le j}\sup_{r\in I}
       |r|^{i+1}|\partial_r^iQ(r)|,\notag\\
 \|(f,g)\|_{\mathsf Y_\varepsilon^s(I)}
 &=\varepsilon^{-1}\max_{0\le j\le s}\sup_{r\in I}|r|^j
       \bigl\{\|\partial_r^jf(r)\|_N
                    +|\partial_r^jg(r)|\bigr\}.    \label{eq:principal-green-spaces}
\end{align}
The factor \(|r|^{-1}\) at order zero in the
\(\mathsf P_\varepsilon^s\)-norm is intentional.

There are linear slow Green operators
\(\mathcal G^{a/r}_{N,\delta}\) such that, writing
\((Z_{\rm sl},Q_{\rm sl})=\mathcal G^{a/r}_{N,\delta}(f,g)\) and
\(\mathfrak p_{\rm sl}=\varepsilon Q_{\rm sl}+\ell_NZ_{\rm sl}\),
\begin{equation}
 \|Z_{\rm sl}\|_{\mathsf Z_\varepsilon^s}
 +\|\mathfrak p_{\rm sl}\|_{\mathsf P_\varepsilon^s}
 +\|Q_{\rm sl}\|_{\mathsf Q_\varepsilon^{s-1}}
 \le C_s\|(f,g)\|_{\mathsf Y_\varepsilon^s}.       \label{eq:principal-green-slow-bound}
\end{equation}
If \((f,g)\in\mathsf Y_\varepsilon^{s+1}\) and the coefficient bounds are
retained through order \(s+2\), then
\begin{equation}
 \|Q_{\rm sl}\|_{\mathsf Q_\varepsilon^s}
 \le C_s\|(f,g)\|_{\mathsf Y_\varepsilon^{s+1}}.  \label{eq:principal-green-extra-q-derivative}
\end{equation}

The endpoint trace maps are
\begin{align}
 \Gamma_a(Z,Q)
 &=\bigl(Z(r_{\rm out}),
   \varepsilon Q(r_{\rm out})
       +\ell_N(r_{\rm out})Z(r_{\rm out})\bigr),\notag\\
 \Gamma_r(Z,Q)
 &=\bigl(Z(-\rho_\delta),
   \varepsilon Q(-r_{\rm out})
       +\ell_N(-r_{\rm out})Z(-r_{\rm out})\bigr).
                                                               \label{eq:principal-green-traces}
\end{align}
There are homogeneous boundary lifts
\begin{equation}
 \mathcal B_a,\mathcal B_r:E_N\times\R
          \longrightarrow\Ker\mathcal L_{N,\delta},
 \qquad \Gamma_a\mathcal B_a=\Gamma_r\mathcal B_r=\Id.       \label{eq:principal-green-lifts}
\end{equation}
Consequently the unique solution with endpoint datum \(b\) is
\begin{equation}
 \binom ZQ
 =\mathcal G^{a/r}_{N,\delta}(f,g)
  +\mathcal B_{a/r}\left[
      b-\Gamma_{a/r}\mathcal G^{a/r}_{N,\delta}(f,g)\right]. \label{eq:principal-green-augmented-inverse}
\end{equation}

The boundary lifts are estimated in \((Z,\mathfrak p)\), not in the top
slow \(Q\)-norm.  Let \(r_a(T)\) solve
\(\dot r=q_0(r)\), \(r_a(0)=r_{\rm out}\), and let \(r_r(T)\) solve
\(\dot r=q_0(r)\), \(r_r(0)=-\rho_\delta\), reaching
\(-r_{\rm out}\) at \(T=T_r\).  Put
\begin{align}
 \Phi_a(T)&=\varepsilon^{-1}\int_0^T
        |w_{0,N}'(r_a(\sigma))|\,\dd\sigma,\notag\\
 \Phi_s(T)&=\varepsilon^{-1}\int_0^T
        w_{0,N}'(r_r(\sigma))\,\dd\sigma,\notag\\
 \Phi_u(T)&=\varepsilon^{-1}\int_T^{T_r}
        w_{0,N}'(r_r(\sigma))\,\dd\sigma.          \label{eq:principal-green-actions}
\end{align}
If \((Z_b,\mathfrak p_b)\) is a boundary lift, then, for
\(0\le j\le s\),
\begin{align}
 \varepsilon^j\bigl\{
  \|\partial_T^jZ_b(T)\|_N
       +|\partial_T^j\mathfrak p_b(T)|\bigr\}
 &\le C_se^{-c\Phi_a(T)}(\|b_z\|_N+|b_p|)          \label{eq:principal-green-attracting-layer}
\end{align}
on the attracting branch.  On the repelling branch write the trace datum
as \(b=(b_s,b_u)\in E_N\times\R\); then
\begin{align}
 \varepsilon^j\bigl\{
  \|\partial_T^jZ_b(T)\|_N
       +|\partial_T^j\mathfrak p_b(T)|\bigr\}
 \le C_s\bigl\{
 e^{-c\Phi_s(T)}\|b_s\|_N
 +e^{-c\Phi_u(T)}|b_u|\bigr\}.                     \label{eq:principal-green-repelling-layer}
\end{align}
Moreover,
\begin{align}
 \Phi_a(T)&\asymp
     \frac{r_{\rm out}^2-r_a(T)^2}{\varepsilon},&
 \Phi_s(T)&\asymp
     \frac{r_r(T)^2-\rho_\delta^2}{\varepsilon},&
 \Phi_u(T)&\asymp
     \frac{r_{\rm out}^2-r_r(T)^2}{\varepsilon},  \label{eq:principal-green-action-comparison}
\end{align}
with dimension-uniform comparison constants.
\end{proposition}

\begin{remark}[No uniform gain of a full slow derivative]
\label[remark]{rem:principal-green-no-derivative-gain}
The endpoint theorem deliberately does not assert a uniform map
\(\mathsf Y_\varepsilon^{s-1}\to
\mathsf Z_\varepsilon^s\times\mathsf Q_\varepsilon^s\).
That assertion is false even for constant coefficients.  On a fixed interior
interval consider
\[
 -\varepsilon Z'+Z=f_\varepsilon,qquad
 h_\varepsilon=\varepsilon^2,qquad
 f_\varepsilon(r)=\varepsilon h_\varepsilon^{s-1}
                         e^{\mathrm i r/h_\varepsilon}.
\]
The \(\mathsf Y_\varepsilon^{s-1}\)-norm of the forcing is bounded, but the
bounded particular solution has
\(|\partial_r^sZ|\asymp1\), so its slow order-\(s\) norm is
\(\asymp\varepsilon^{-1}\).

The one-order loss in \(Q\) is also genuine.  Freeze
\(B=-1\), \(q_0=-1\), \(z_0'=0\), \(\ell=r_*>0\), and
\(w_0'=-r_*\), and take
\(f_\varepsilon=\varepsilon^{s+1}e^{\mathrm i r/\varepsilon}\),
\(g=0\).  The particular solution has
\(|Z|\asymp\varepsilon^{s+1}\) and
\(|Q|\asymp r_*\varepsilon^s\).  Its slow \(Q\)-norm is bounded through
order \(s-1\), and its input \(\mathsf Y_\varepsilon^s\)-norm is bounded,
whereas the order-\(s\) \(Q\)-norm grows like
\(c\varepsilon^{-1}\).  Taking real or imaginary parts gives the same
examples in the real Banach spaces.  They use the exact traces of the
particular solutions and therefore are not endpoint boundary layers.
Thus the coefficient hypotheses in the proposition alone cannot yield a
same-order estimate for \(Q\).
\end{remark}

The principal cancellation coordinate has an exact nonlinear counterpart.
It is the resource defect itself, rather than a formal change of variables.
For a nonvanishing speed \(q\), define
\begin{align}
 \mathcal D_N(r,z,q,\eta)
 =K\sum_{k=0}^mL_{k,N}(\eta)
 \bigl[\{r-r_k^q(r)\}\mathbf1
       +z(r)-z(r_k^q(r))\bigr],                   \label{eq:nonlinear-tracker-delay}\\
 \mathcal D_{c,N}&=\pi_N^T\mathcal D_N,
 &\mathcal D_{\perp,N}&=P_{\perp,N}\mathcal D_N,\notag
\end{align}
where \(r_k^q=\Phi_q^{-\delta\theta_k}\).  Put
\begin{equation}
 F_N(r,z)=\pi_N^T\mathcal N_N(r\mathbf1+z).        \label{eq:nonlinear-tracker-F}
\end{equation}

\begin{proposition}[Exact nonlinear resource-defect normal form]
\label[proposition]{prop:nonlinear-resource-defect-normal-form}
Retain the shared-resource equation and the intervals of
\cref{prop:principal-tracker-endpoint-green}.  Let
\(z=z_{0,N}+Z\), \(q=q_0+Q\), and suppose all backtracks remain in one
common collar.  Set
\begin{align}
 \mathcal R_N^{(2)}(r,Z)
 &={\mathcal N}_N(x_{0,N}+Z)-{\mathcal N}_N(x_{0,N})
      -D{\mathcal N}_N(x_{0,N})Z,\notag\\
 h_{c,N}&=\pi_N^T\mathcal R_N^{(2)},\qquad
 h_{\perp,N}=P_{\perp,N}\mathcal R_N^{(2)},\notag\\
 \mathcal D_{0,N}&=\mathcal D_N(r,z_{0,N},q_0,\eta),\notag\\
 \Delta\mathcal D_N
 &=\mathcal D_N(r,z_{0,N}+Z,q_0+Q,\eta)-\mathcal D_{0,N},\notag\\
 R_{z,N}&=\varepsilon
   \{q_0z_{0,N}'-\mathcal D_{0,\perp,N}\},\notag\\
 R_{q,N}&=\varepsilon q_0
   (\mathcal D_{0,c,N}-q_0)',\notag\\
 \mathcal C_N(r)&=B_N(r)+z_{0,N}'(r)\ell_N(r).   \label{eq:nonlinear-tracker-data}
\end{align}
Define the seed resource and the nonlinear resource defect by
\begin{align}
 w_{\rm seed,N}
 &=w_{0,N}+\varepsilon\mathcal D_{0,c,N}-\varepsilon q_0,\notag\\
 \mathfrak p_{\rm nl}
 &=w_{\rm seed,N}-w\notag\\
 &=\varepsilon Q+\ell_NZ+h_{c,N}
       -\varepsilon\Delta\mathcal D_{c,N}.        \label{eq:nonlinear-resource-defect}
\end{align}
Then the two projected identities in
\cref{eq:exact-outer-voltage-invariance,eq:exact-outer-resource-invariance}
are equivalent to
\begin{align}
 \varepsilon qZ'
 ={}&\mathcal C_NZ-z_{0,N}'\mathfrak p_{\rm nl}-R_{z,N}
       +z_{0,N}'h_{c,N}-h_{\perp,N}\notag\\
 &+\varepsilon\{
       \Delta\mathcal D_{\perp,N}
       -z_{0,N}'\Delta\mathcal D_{c,N}\},         \label{eq:nonlinear-tracker-Z-normal-form}\\
 \varepsilon q\mathfrak p_{\rm nl}'
 ={}&w_{0,N}'\mathfrak p_{\rm nl}-w_{0,N}'\ell_NZ
       -w_{0,N}'h_{c,N}\notag\\
 &+\varepsilon w_{0,N}'\Delta\mathcal D_{c,N}
       +\varepsilon^2q(\mathcal D_{0,c,N}-q_0)',   \label{eq:nonlinear-tracker-p-normal-form}\\
 \varepsilon Q
 ={}&\mathfrak p_{\rm nl}-\ell_NZ-h_{c,N}
       +\varepsilon\Delta\mathcal D_{c,N},
 &q={}&q_0+Q,                                      \label{eq:nonlinear-tracker-Q-reconstruction}\\
 w={}&w_{\rm seed,N}-\mathfrak p_{\rm nl}.        \label{eq:nonlinear-tracker-w-reconstruction}
\end{align}
In particular, no derivative of \(h_{c,N}\),
\(\Delta\mathcal D_{c,N}\), or \(Q\) occurs in the nonlinear forcing.
At the origin the right-hand sides have affine value
\((-R_{z,N},\varepsilon R_{q,N})\).  If
\[
 \mathscr K_{\perp/c,N}[\widehat Z,\widehat Q]
 =D\Delta\mathcal D_{\perp/c,N}(0,0)
       [\widehat Z,\widehat Q],
\]
then the exact differential is
\begin{align}
 \varepsilon q_0\widehat Z'
 &=\mathcal C_N\widehat Z-z_{0,N}'\widehat p
   +\varepsilon(\mathscr K_{\perp,N}
                  -z_{0,N}'\mathscr K_{c,N}),\notag\\
 \varepsilon q_0\widehat p'
 &=w_{0,N}'\widehat p-w_{0,N}'\ell_N\widehat Z
   +\varepsilon w_{0,N}'\mathscr K_{c,N}
   +\varepsilon^2(\mathcal D_{0,c,N}-q_0)'\widehat Q,\notag\\
 \varepsilon\widehat Q
 &=\widehat p-\ell_N\widehat Z
   +\varepsilon\mathscr K_{c,N}.                 \label{eq:nonlinear-tracker-full-differential}
\end{align}
Thus \cref{eq:principal-green-transformed-system} is the retained leading
differential block; the full linearization additionally contains the
displayed \(\varepsilon\)-weighted delay/backtrack terms.

Around each admissible base point, after shrinking the product neighborhood,
the algebraic relation
\cref{eq:nonlinear-tracker-Q-reconstruction} locally determines \(Q\) from
\((Z,\mathfrak p_{\rm nl})\) in \(C^0\).  More precisely, on every common
flow tube on which \(-q\) is bounded above and below and
\(\|\mathbf1+z_{0,N}'+Z'\|_N\le M_X\),
\begin{equation}
 \|D_Q\Delta\mathcal D_{c,N}[H]\|_{C^0}
 \le C\delta M_X\|H\|_{C^0}.                     \label{eq:nonlinear-tracker-Q-contraction}
\end{equation}
Hence the coordinate change is one-to-one for small \(\delta\), uniformly
in \(N\).  The reduced structural directions have no direct contribution to
\(\mathcal D_{c,N}\), because \(\pi_N^TE_k=0\).
\end{proposition}

The remaining endpoint object is a genuine RFDE history, not an independent
function \(q_{\rm old}\).  The next result constructs its strong compatibility
jets without adding rows to the Green boundary operator.  In fold time write
the raw fields as
\begin{align}
 \mathfrak F_{N,\delta,\nu,\eta}(\phi,\omega)
 ={}&\delta^{-1}\bigl[
  (2/3-\omega)\mathbf1
  -\mathcal N_N(\phi(0)-\mathbf1)
  +\mathcal A_N\{\phi(0)-\mathbf1\}\bigr]\notag\\
 &+\delta K\sum_{k=0}^mL_{k,N}(\eta)
       \{\phi(0)-\phi(-\theta_k)\},\notag\\
 \mathfrak G_{N,\delta,\nu}(\phi)
 ={}&\delta\{\pi_N^T\phi(0)-1-\delta^2\nu\}.       \label{eq:raw-fold-time-fields}
\end{align}

\begin{proposition}[Compatible endpoint-jet collar chart]
\label[proposition]{prop:compatible-endpoint-jet-collar}
Fix an integer \(h\ge3\).  Suppose at least one delay is positive, put
\(\theta_+=\min\{\theta_k:\theta_k>0\}\), and choose
\(c_H>0\) and small \(\delta\) so that \(4c_H\delta<\theta_+\).
Decrease \(r_{\rm out}\le1\) if necessary.
For an endpoint \(e\) satisfying
\(\rho_\delta\le|e|\le r_{\rm out}\), set
\(\kappa_e=\kappa_{N,\delta}(e)\) and define
\begin{equation}
 \|\psi\|_{\mathrm{cap},h,e}
 =\max_{0\le j\le h}\delta^j
   \sup_{-\theta_m\le\theta\le0}
    e^{2\kappa_e\theta}\|\partial_\theta^j\psi(\theta)\|_N.   \label{eq:strong-collar-norm}
\end{equation}
There is a dimension-uniform linear Hermite operator
\begin{equation}
 \mathcal E_{\delta,h}(a_0,\ldots,a_h)(\theta)
 =\sum_{j=0}^h\delta^ja_j\chi_j(\theta/\delta),    \label{eq:collar-Hermite-lift}
\end{equation}
where the scalar one-sided \(C^\infty(({-}\infty,0])\) functions \(\chi_j\)
are supported, relative to that half-line, in \([-c_H,0]\) and satisfy
\(\chi_j^{(i)}(0)=\delta_{ij}\).  It obeys
\begin{equation}
 \|\mathcal E_{\delta,h}a\|_{\mathrm{cap},h,e}
 \le C_h\max_{0\le j\le h}\delta^j\|a_j\|_N.      \label{eq:collar-Hermite-bound}
\end{equation}
Put
\begin{equation}
 \mathscr K_{N,\delta}^h
 =\{\widehat\gamma\in C^h([-\theta_m,0];\R^N):
        \partial_\theta^j\widehat\gamma(0)=0,\ 0\le j\le h\}.
                                                               \label{eq:collar-core-space}
\end{equation}

On every fixed bounded data ball whose endpoint values lie in a sufficiently
small weighted affine neighborhood of a reference critical collar, every
\begin{equation}
 (Z_e,p_e,\widehat\gamma;\nu,\eta)
 \in E_N\times\R\times\mathscr K_{N,\delta}^h
       \times\R\times\mathfrak T_N^{\rm red}       \label{eq:collar-chart-data}
\end{equation}
determines a unique history--resource pair
\begin{equation}
 \Gamma_e^h(Z_e,p_e,\widehat\gamma;\nu,\eta)
   =(\phi_e,\omega_e)                               \label{eq:collar-chart}
\end{equation}
of the form \(\phi_e=\widehat\gamma+\mathcal E_{\delta,h}a\), satisfying
\begin{align}
 \phi_e(0)&=V_{0,N}(e)+Z_e,&
 \omega_e&=w_{\rm seed,N}(e)-p_e,                 \label{eq:collar-current-data}\\
 \partial_\theta^j\phi_e(0)
 &=\mathfrak V_{j,N}(\phi_e,\omega_e;\delta,\nu,\eta),
 &1&\le j\le h.                                   \label{eq:collar-compatibility}
\end{align}
Here \(\mathfrak V_{1,N}=\mathfrak F_{N,\delta,\nu,\eta}\), and
\(\mathfrak V_{j,N}\) is obtained by differentiating the two raw equations
\(j-1\) times, using \(\mathfrak G\) for every resource derivative.
Thus \cref{eq:collar-compatibility} is precisely the order-\(h\) raw-RFDE
compatibility condition.  The recursion for the endpoint jets is triangular:
the right-hand side at order \(j\) uses only endpoint jets of order below
\(j\) and the prescribed core jets at the positive delays.

For fixed parameters, two chart values satisfy the dimension-uniform tame
estimate
\begin{align}
 \|\phi_e-\widetilde\phi_e\|_{\mathrm{cap},h,e}
 +\frac{|\omega_e-\widetilde\omega_e|}{|e|}
 \le C_h\biggl\{
  \|Z_e-\widetilde Z_e\|_N
  +\frac{|p_e-\widetilde p_e|}{|e|}
  +\|\widehat\gamma-\widetilde{\widehat\gamma}\|_{
        \mathrm{cap},h,e}\biggr\}.                \label{eq:collar-chart-tame-bound}
\end{align}
The chart, as a map into the fixed \(\theta\)-history space, is three times
Fr\'{e}chet differentiable in all its arguments, with analogous triangular
tame bounds, uniformly in \(N\).  In particular it supports every
rectangular parameter jet in
\((C_\nu^1C_\eta^2)\cap(C_\nu^2C_\eta^1)\).

If the collective collar coordinate
\begin{equation}
 r_{\phi_e}(\theta)=\pi_N^T\{\phi_e(\theta)-\mathbf1\}         \label{eq:collar-collective-coordinate}
\end{equation}
is monotone and
\(-q_+\le\delta^{-1}r_{\phi_e}'\le-q_-\), its image carries the unique
old-speed function
\begin{equation}
 q_{\rm old}(r_{\phi_e}(\theta))
 =\delta^{-1}r_{\phi_e}'(\theta).                  \label{eq:collar-derived-old-speed}
\end{equation}
It belongs to \(C^{h-1}\), includes its complete endpoint jet, and is not
an additional boundary datum.  More precisely,
\begin{equation}
 r_{\phi_e}(\theta)=\Phi_{q_{\rm old}}^{\delta\theta}(e),
 \qquad r_{\phi_e}(0)=e,                           \label{eq:collar-backtrack-bridge}
\end{equation}
and, with
\[
 \mathcal D_{c,N}[\phi_e]
 =K\sum_{k=0}^m\pi_N^TL_{k,N}(\eta)
       \{\phi_e(0)-\phi_e(-\theta_k)\},
\]
the first compatibility row gives the exact endpoint bridge
\begin{equation}
 \varepsilon\{q_{\rm old}(e)-q_0(e)\}
 =p_e-\ell_N(e)Z_e-h_{c,N}(e,Z_e)
  +\varepsilon\{\mathcal D_{c,N}[\phi_e]
                    -\mathcal D_{0,c,N}(e)\}.     \label{eq:collar-speed-defect-bridge}
\end{equation}
Thus the right-hand side is exactly the reconstruction of
\(\varepsilon Q_e\) in
\cref{eq:nonlinear-tracker-Q-reconstruction}.  The exact nonlinear scalar
trace is simply
\begin{equation}
 \mathfrak p_{\rm nl}(e)=w_{\rm seed,N}(e)-\omega_e=p_e.      \label{eq:collar-nonlinear-trace}
\end{equation}
Consequently the finite Green trace count remains
\((N-1)+1=N\).  On the repelling branch the inner value \(p_e\) may be the
output of the mixed BVP; evaluating \(\Gamma_e^h\) at that value imposes no
second scalar boundary row.  The core variable \(\widehat\gamma\) is a
collar gauge, not an additional Fredholm row.

The pullback
\(q_{\rm old}\circ r_{\phi_e}=\delta^{-1}r_{\phi_e}'\) carries the same
three parameter jets in the fixed \(\theta\)-chart.  No lossless Eulerian
fixed-\(r\) assertion is made: on a common image interval, a parameter
derivative of total order \(0\le b\le\min\{3,h-1\}\) generally has only
\(D_\lambda^bq_{\rm old}\in C^{h-1-b}\); a full third Eulerian parameter
jet therefore requires \(h\ge4\).  Nor is
\(\Gamma_e^h\) a compatible right inverse for arbitrary Green sources; it
is a chart of raw-compatible histories, and matching it to such a right
inverse is a separate hybrid step.  If there is no positive delay, the
statement reduces to the corresponding finite-dimensional jet recursion.
\end{proposition}

The strong collar chart does not make every endpoint Green lift slow.  The
following scalar mode identifies the remaining hybrid boundary-layer gate.

\begin{proposition}[Finite endpoint scale does not control the slow speed]
\label[proposition]{prop:endpoint-scale-not-slow-speed}
Let \(I=[\rho,R]\) with fixed \(R>0\), take
\(q_0=-c\), \(w_0'(r)=-r/c\), \(\ell=0\), and \(Z=0\), where \(c>0\).
For every fixed \(0<A<c/2\), the homogeneous scalar equation has the solution
\begin{equation}
 Q_A(r)=A\exp\left\{\frac{r^2-R^2}{2c^2\varepsilon}\right\},
 \qquad \mathfrak p_A(r)=\varepsilon Q_A(r).       \label{eq:speed-boundary-layer}
\end{equation}
It retains the strict speed sign \(q_0+Q_A\le-c/2\) and has endpoint datum
\(\mathfrak p_A(R)=\varepsilon A\), hence the finite endpoint scale is
\begin{equation}
 \varepsilon^{-1}R^{-1}|\mathfrak p_A(R)|=A/R.    \label{eq:speed-boundary-layer-trace-scale}
\end{equation}
Nevertheless,
\begin{equation}
 \|Q_A\|_{\mathsf Q_\varepsilon^0}
 \ge\frac{RA}{\varepsilon},
 \qquad R|Q_A'(R)|=\frac{R^2A}{c^2\varepsilon},    \label{eq:speed-boundary-layer-blowup}
\end{equation}
so it does not belong to the uniform slow-speed class of
\cref{prop:physical-backtrack-calculus}.  If \(w'=r/(q_0+Q_A)\), then
\begin{equation}
 |w'(R)-w_0'(R)|=\frac{RA}{c(c-A)},                \label{eq:speed-boundary-layer-resource}
\end{equation}
which is not \(O(\varepsilon)\).  Thus the finite endpoint scale alone does
not supply the slow-speed hypothesis.  Coupling the collar chart to the
finite endpoint theorem still requires a slow/action hybrid backtrack
estimate, or an explicit first-delay buffer, before it can close the
nonlinear tracker.
\end{proposition}

There is nevertheless an intrinsic quotient once an actual physical history
has been constructed.  The point is to quotient by the transported tangent
line, rather than by an eigenvector of a frozen voltage--resource matrix; a
fixed complement is only a coordinate gauge for that quotient.

\begin{proposition}[Stationary-coordinate quotient along a physical history]
\label[proposition]{prop:physical-history-phase-quotient}
Fix \(N,\delta,\nu,\eta\), and in the raw equation replace
\(B_{k,N}(\zeta)\) by the structural-ball layer \(L_{k,N}(\eta)\).
Let \((\bar v(t),\bar w(t))\) be a classical solution of this physical
equation and rescale it to fold time by
\[
 v_*(s)=\bar v(s/\delta),\qquad w_*(s)=\bar w(s/\delta).
\]
Thus the fold-time delays are \(\theta_k\).  Let
\(Z_*(s)=(v_{*,s},w_*(s))\in\mathscr Y_N\) be compatible and \(C^2\) on an
interval enlarged by one delay, and let
\[
 r_*(s)=\pi_N^T\{v_*(s)-\mathbf1\}.
\]
Suppose \(r_*'(s)\ne0\) throughout that interval, and denote the variational
solution process of the rescaled RFDE along \(Z_*\) by
\(\mathcal U_*(s,u)\).  Then
\begin{equation}
 \tau_*(s)=\frac1{r_*'(s)}
 \left(\theta\longmapsto v_*'(s+\theta),\;w_*'(s)\right)
 \in\mathscr Y_N                                      \label{eq:physical-normalized-tangent}
\end{equation}
is the unique vector in the time-tangent line satisfying
\(\varrho_N\tau_*(s)=1\).  Its line is transported exactly:
\begin{equation}
 \mathcal U_*(s,u)\tau_*(u)
 =\frac{r_*'(s)}{r_*'(u)}\tau_*(s).                  \label{eq:physical-tangent-transport}
\end{equation}
Consequently
\begin{equation}
 Q_*(s)=\Id-\tau_*(s)\varrho_N,
 \qquad
 \widehat{\mathcal U}_*(s,u)
 =Q_*(s)\mathcal U_*(s,u)|_{\ker\varrho_N}          \label{eq:physical-quotient-process}
\end{equation}
defines an exact linear process on the fixed stationary-coordinate gauge
section
\(\ker\varrho_N\): for \(u\le t\le s\),
\begin{equation}
 \widehat{\mathcal U}_*(s,t)
 \widehat{\mathcal U}_*(t,u)
 =\widehat{\mathcal U}_*(s,u).                      \label{eq:physical-quotient-cocycle}
\end{equation}
Equivalently, the intrinsic object is the variational process induced on
\(\mathscr Y_N/\operatorname{span}\{\tau_*(s)\}\), written in the physical
stationary-coordinate gauge.  Another bounded transversal gives a conjugate
representation of the same quotient process provided it is uniformly
transverse to the tangent line and its projections, quotient
identifications, and inverse identifications have common bounds.  In every
history norm for which
\(\varrho_N\) is bounded,
\begin{equation}
 \|Q_*(s)\|\le1+\|\tau_*(s)\|\,\|\varrho_N\|.       \label{eq:physical-quotient-bound}
\end{equation}
The construction is unchanged if the time tangent is multiplied by a
nonzero scalar before imposing \(\varrho_N\tau=1\); no eigenvector or
spectral projection of a frozen full system enters it.
\end{proposition}

The row \(\varrho_N\) need not be dual-covariant, and
\(\ker\varrho_N\) need not be invariant; the projection at every step in
\cref{eq:physical-quotient-process} is essential.  A dual-covariant phase
functional annihilating the stable and strong-unstable bundles could be
defined only after the missing normal dichotomy has been proved.  Thus the
proposition constructs neither such a phase functional nor an invariant
normal bundle.

The monotonicity in \cref{eq:outer-speed-sign} is therefore the correct
leading input for phase removal, but it is not yet the hypothesis of
\cref{prop:physical-history-phase-quotient}: one must first construct a true
positive-\(\delta\) history with \(r_*'\ne0\), and then prove dimension-uniform
bounds and the required parameter jets for
\cref{eq:physical-normalized-tangent,eq:physical-quotient-process}.  In
particular, the fixed row in
\cref{prop:backward-asymptotic-nonselection} can normalize the tangent of
each already selected member, but the row alone selects neither that physical
history nor its transported tangent line.  The present proposition is not a
global anchor or history-selection theorem.

\begin{proposition}[Exact resource-gauge normal equation]
\label[proposition]{prop:exact-resource-gauge-quotient}
Retain \cref{prop:exact-outer-history-equation} for the shared-resource law,
and let \(r_s\) be a generated orbit in fold time:
\[
 r_s'=\delta q(r_s),\qquad
 Z_*(s)=\mathcal H_{V,w,q}(r_s).
\]
Assume \(w'(r_s)\ne0\) on a delay-enlarged orbit interval, and denote the
variational process along \(Z_*\) by \(\mathcal U_*(s,\sigma)\).  The
automatic \(C^2\) regularity in
\cref{prop:exact-outer-history-equation} applies.  Write
\begin{align}
 r_{s+\theta}&=\Phi_q^{\delta\theta}(r_s),&
 \chi_s(\theta)&=\frac{q(r_{s+\theta})}{q(r_s)},\notag\\
 e_s(\theta)&=\chi_s(\theta)V'(r_{s+\theta}),&
 E_s^c&=(e_s,w'(r_s)).                              \label{eq:resource-gauge-center-history}
\end{align}
Then \(E_s^c=\partial_r\mathcal H_{V,w,q}(r_s)\), and
\begin{equation}
 \mathcal U_*(s,\sigma)E_\sigma^c
 =\frac{q(r_s)}{q(r_\sigma)}E_s^c.                 \label{eq:resource-gauge-center-cocycle}
\end{equation}
In particular, the genuine time-translation tangent is
\(\delta q(r_s)E_s^c\), not a frozen slow eigenvector.

Let \((\xi_s,\omega(s))\) solve the full variational equation along \(Z_*\).
Define
\begin{equation}
 a(s)=\frac{\omega(s)}{w'(r_s)},\qquad
 \psi_s(\theta)=\xi(s+\theta)-e_s(\theta)a(s),       \label{eq:resource-gauge-coordinates}
\end{equation}
and put
\begin{equation}
 h_s(\theta)=\frac{e_s(\theta)}{w'(r_s)},\qquad
 H_s=h_s(0)=\frac{V'(r_s)}{w'(r_s)}.                \label{eq:resource-gauge-h}
\end{equation}
At each fixed \(s\), the history \(\psi_s\) is unchanged when any multiple of
\(E_s^c\) is added to the full variation.  With \(u(s)=\psi_s(0)\), the exact triangular
quotient equations are
\begin{align}
 a'(s)&=\delta q'(r_s)a(s)
       +\frac{\delta}{w'(r_s)}\pi_N^Tu(s),           \label{eq:resource-gauge-a}\\
 (\partial_s-\partial_\theta)\psi_s(\theta)
 &=-\delta h_s(\theta)\pi_N^Tu(s),                  \label{eq:resource-gauge-transport}\\
 u'(s)&=\delta^{-1}\mathcal J_*(s)u(s)
 +\delta K\sum_{k=0}^mL_{k,N}(\eta)
       \{u(s)-\psi_s(-\theta_k)\}
 -\delta H_s\pi_N^Tu(s),                            \label{eq:resource-gauge-boundary}
\end{align}
where the transport identity is for
\(-\theta_m\le\theta<0\), in the mild characteristic sense for continuous
histories; \(\theta=0\) is governed by
\cref{eq:resource-gauge-boundary}.  Here
\begin{equation}
 \mathcal J_*(s)=\mathcal A_N-\diag\!\left(
 2c_N\circ\{V(r_s)-\mathbf1\}
 +\beta\{V(r_s)-\mathbf1\}^{\circ2}\right).         \label{eq:resource-gauge-current-matrix}
\end{equation}

This transport--boundary system has an exact realization on an ordinary
voltage history.  For \(u_s(\sigma)=u(s+\sigma)\), define the Volterra
operator
\begin{equation}
 (\mathcal K_s\phi)(\theta)
 =\delta\int_\theta^0
 h_{s+\sigma}(\theta-\sigma)\pi_N^T\phi(\sigma)\,\dd\sigma.
                                                               \label{eq:resource-gauge-volterra}
\end{equation}
Then
\begin{equation}
 \psi_s=(\Id-\mathcal K_s)u_s,                       \label{eq:resource-gauge-history-conjugacy}
\end{equation}
and \(u\) satisfies the standard nonautonomous RFDE
\begin{align}
 u'(s)={}&\delta^{-1}\mathcal J_*(s)u(s)
 +\delta K\sum_{k=0}^mL_{k,N}(\eta)
 \bigl[u(s)-u(s-\theta_k)
       +(\mathcal K_su_s)(-\theta_k)\bigr]\notag\\
 &-\delta H_s\pi_N^Tu(s).                           \label{eq:resource-gauge-standard-rfde}
\end{align}

Suppose, on the whole delay collar, that \(q\), its reciprocal, and \(V'\)
have common bounds, and that
\begin{equation}
 |w'(r)|\ge c_w|r|,\qquad |r|\ge c_\rho\rho_\delta. \label{eq:resource-gauge-collar-bounds}
\end{equation}
For every \(\kappa_*\ge0\), use
\[
 \|\phi\|^\sharp_{\kappa_*}
 =\sup_{-\theta_m\le\theta\le0}
 e^{2\kappa_*\theta}\|\phi(\theta)\|_N.
\]
Then, uniformly in \(N\) and the orbit point,
\begin{equation}
 \|\mathcal K_s\|
 \le C\frac{\delta}{\rho_\delta}
 \le\frac{C}{S_\delta}.                             \label{eq:resource-gauge-volterra-bound}
\end{equation}
There is \(\delta_0>0\) such that, for \(0<\delta<\delta_0\),
\begin{equation}
 \|(\Id-\mathcal K_s)^{-1}\|\le2,\qquad
 \|(\Id-\mathcal K_s)^{-1}-\Id\|\le\frac{C}{S_\delta}.
                                                               \label{eq:resource-gauge-volterra-inverse}
\end{equation}
Conversely, on this range
\begin{equation}
 u_s=(\Id-\mathcal K_s)^{-1}\psi_s,                 \label{eq:resource-gauge-history-inverse}
\end{equation}
so \cref{eq:resource-gauge-history-conjugacy} is a two-sided coordinate
equivalence between the quotient transport--boundary system and
\cref{eq:resource-gauge-standard-rfde}; the scalar \(a\) is then recovered
from \cref{eq:resource-gauge-a}.
The full moving norm
\begin{equation}
 \|(\xi_s,\omega)\|_{*,s}
 =\|\xi_s\|^\sharp_{\kappa_*}
   +\frac{|\omega|}{|w'(r_s)|}                      \label{eq:resource-gauge-anisotropic-norm}
\end{equation}
is uniformly equivalent to
\(\|\psi_s\|^\sharp_{\kappa_*}+|a(s)|\).
\end{proposition}

Thus the apparently large term \(-\omega\mathbf1/\delta\) in the raw
fold-time variational equation is removed exactly only after the genuine
tracker tangent is known.  The remaining Volterra coordinate correction is
\(O(S_\delta^{-1})\), but
\cref{prop:exact-resource-gauge-quotient} still assumes the tracker and the
collar bounds; it does not prove them.  Uniform parameter derivatives further
require relative bounds on the jets of \(w'\), not merely absolute bounds.

The proofs, including the dimension-uniform block graph construction, are in
\cref{sec:supplementary-outer-skeleton}.  Three qualifications are essential.
First, \cref{eq:outer-stable-layer-semigroup} concerns the leading
frozen-resource voltage matrix \(\mathcal J^{\rm fv}_{N}(r)\).  Its
positive-\(\delta\) frozen-resource RFDE upgrade is precisely
\cref{prop:frozen-voltage-history-splitting}; neither result includes the
resource tangent or a nonautonomous slow history.  Second, the action is that of the principal
normal equation \(\delta^2n_T=\lambda_{c,N}(r(T))n\), not yet an action
estimate for an invariant RFDE slow-history graph.  Third, at
\(r=\pm\rho_\delta\), so that \(X=r/\delta\), the critical curve has
\begin{equation}
 d_{\rm crit}
 =\frac{2/3-w_{0,N}(r)}{\delta^2}-\alpha X^2
      +\frac1{2\alpha}
 =\frac1{2\alpha}+O_{\rm unif}(\delta S_\delta^3),            \label{eq:outer-critical-fold-offset}
\end{equation}
whereas the retained canard tube has \(|d|<\delta^\kappa\).  Thus, in the
present \(r\)-parameterization and fold chart, a physical outer-history
theorem must at least produce the leading finite-\(\delta\) correction
\begin{equation}
 w^{{\rm slow},\diamond,\pm}_{N,\delta,\nu,\zeta}(\pm\rho_\delta)
 =w_{0,N}(\pm\rho_\delta)+\frac{\delta^2}{2\alpha}
   +E^{\diamond,\pm}_{N,\delta,\nu,\zeta}.          \label{eq:outer-required-resource-correction}
\end{equation}
A sufficient resource-coordinate remainder for entry into
\(|d|<\delta^\kappa\) would be
\(E^{\diamond,\pm}_{N,\delta,\nu,\zeta}
=o(\delta^{2+\kappa})\), uniformly in the displayed indices on a fixed
admissible compact \((\nu,\zeta)\)-window.  Even that
would not prove exact overlap without the corresponding voltage and full
history estimates.

Nor can such a theorem obtain a unique history merely by asking for a
bounded complete backward extension converging in an unnormalized history
norm to a reference branch.
\Cref{prop:backward-asymptotic-nonselection} gives an explicit RFDE
counterexample, even for families which are superalgebraically close in a
strong history norm and share one phase row.  These qualitative conditions
alone prove neither uniqueness nor the required weighted parameter jets.
Any complete-history construction for this RFDE must therefore specify an
additional model-dependent selection or normalization and prove its uniform
jets.
